\documentclass{situationist-preprint}

\title{The Rationality of Nostalgia}
\msid{EXI-14-a}
\msfamily{EXI-14}
\msdiscipline{Meaning, Death and the Absurd}
\mstype{Principal manuscript}
\mskeywords{nostalgia; emotion; fittingness; past; memory; temporal attitudes}

\begin{document}
\maketitle

\begin{abstract}
Is nostalgia rational? It has been diagnosed as a disease, condemned as sentimental distortion and politically suspect, and defended by psychologists as a resource for meaning and social connection. We argue that nostalgia is rational in moderation: it is rational when it tracks real goods of the past. On the account of emotion as evaluative representation, nostalgia represents some part of one's past as good and as no longer available, and it is fitting when the past contained that good and the good is no longer to be had. Its two characteristic failures are distortion, when the remembered good was not there or not as remembered, and excess, when the attitude is out of proportion to the good or crowds out engagement with present goods. Neither failure is intrinsic to nostalgia. The account requires no tensed metaphysics: the goods of the past are real and fixed, and what is lost is the subject's access to them from where she now stands in time. Nostalgia is thus rational in the same way grief is, and for the same reasons.
\end{abstract}

\section{The problem}

The word was coined in 1688 by the physician Johannes Hofer for a disease of Swiss mercenaries pining for home. It later lost its medical sense and acquired a temporal one: a bittersweet longing for one's own past. Its critics regard it as irrational in two ways. It distorts, since it remembers the past selectively and as better than it was; and it disables, since it turns the subject away from the present. \citet{Boym2001} distinguished a restorative nostalgia, which seeks to rebuild a lost home and fuels nationalist myth, from a reflective nostalgia, which dwells on longing and loss without the illusion of return. Psychologists have found that nostalgia is a common and largely positive emotion that raises self-continuity, social connectedness and the sense that life is meaningful \citep{WildschutEtAl2006, SedikidesEtAl2008}. \citet{Howard2012} gave a philosophical analysis of nostalgia on which its characteristic distortions are not essential to it.

These findings and arguments concern whether nostalgia is good for us. Our question is whether it is rational, and we take the question in the sense in which an emotion is rational when it fits its object.

\section{Framework}

\begin{principle}[Emotions as evaluations]\label{p:emo}
An emotion represents its object as having an evaluative property, and it is \emph{fitting} when the object has that property (MND-17; \citealp{DArmsJacobson2000}).
\end{principle}

\begin{definition}[Nostalgia]\label{d:nost}
\emph{Nostalgia} is a bittersweet emotion directed at a part of one's own past, which represents it as having contained a good, and that good as no longer available to the subject.
\end{definition}

\begin{definition}[Rational nostalgia]\label{d:rat}
An episode of nostalgia is \emph{rational} if it is fitting (the past contained the good, and the good is no longer available), and proportionate (its intensity and its role in the subject's attention correspond to the size of the good and do not displace response to present goods).
\end{definition}

Proportion is a general condition on the rationality of emotion \citep{deSousa1987}: grief that is fitting may still be excessive, and fear of a real but slight danger may be out of proportion to it.

\section{The argument}

\begin{lemma}\label{l:fit}
Nostalgia is fitting when it tracks a real good of the past that is no longer available.
\end{lemma}

\begin{proof}
By Definition~\ref{d:nost}, nostalgia represents a past as containing a good that is no longer available. By Principle~\ref{p:emo}, it is fitting when the past did contain the good and it is no longer available. These are the conditions stated.
\end{proof}

\begin{theorem}\label{t:main}
Nostalgia is rational in moderation: an episode is rational if and only if it tracks a real good of the past that is no longer available and is proportionate to that good.
\end{theorem}

\begin{proof}
By Definition~\ref{d:rat}, rationality is fittingness together with proportion. By Lemma~\ref{l:fit}, fittingness is the tracking of a real past good no longer available. Proportion is moderation. The conjunction is the stated condition.
\end{proof}

The two familiar charges against nostalgia correspond to the two conditions. Distortion is a failure of fit: the remembered summer was not as golden, the old neighborhood was not as harmonious, the nation never was as it is now imagined. Excess is a failure of proportion: the longing occupies more of one's attention than the good warrants, or forecloses goods available now. Since neither failure is entailed by Definition~\ref{d:nost}, neither is intrinsic to nostalgia. Boym's restorative nostalgia typically fails both conditions, since it constructs a past that did not exist and seeks to return to it; her reflective nostalgia, which knows the past is gone and does not falsify it, is the form that the theorem allows.

The bitter and the sweet in nostalgia correspond to its two representational components: the sweetness to the good, the bitterness to its unavailability. This is why nostalgia is fitting in the same way as grief and gratitude: grief responds to a good that was real and is lost, gratitude to a good that was received.

\section{Objections and replies}

\begin{objection}[The past is gone]\label{o:gone}
Nostalgia represents its object as lost. If the past exists and is fixed, nothing is lost, and nostalgia misrepresents.
\end{objection}

\begin{reply}
TSM-03 argues that the past exists, elsewhere in time, as distant places exist elsewhere in space. That supports, rather than undermines, the first component of nostalgia: the good it remembers is real and nothing can now make it not to have been. What nostalgia represents as lost is not the good but the subject's access to it. The good is located at a time earlier than the subject's present, and she cannot be there again, as the emigrant cannot be in the home country while abroad, with the difference that the temporal distance cannot be crossed. The relation that grounds the loss is the subject's position in time relative to the good, and it holds on a tenseless view as much as on any other. Nostalgia is in this respect like homesickness, from which it historically derives.
\end{reply}

\begin{objection}[Selective memory]\label{o:select}
Memory is reconstructive and biased toward the positive. Nostalgia always tracks a reconstructed past, not the real one.
\end{objection}

\begin{reply}
Selective attention to the good in a past is not distortion, if the good was there. A person who recalls the warmth of a family meal without recalling the quarrel that preceded it remembers something real. Distortion occurs when what is recalled was not there or was not as good as it is recalled to be, and the subject is answerable to this: nostalgia, like other emotions, is defeated by evidence that its object lacks the property it represents. That memory is fallible shows only that nostalgia can be unfitting, which the theorem allows.
\end{reply}

\begin{objection}[Uselessness]\label{o:useless}
Nothing can be done about the past. An emotion that motivates no possible action is irrational.
\end{objection}

\begin{reply}
Fittingness does not require that an emotion prompt action, as grief and admiration show. Nostalgia does, moreover, have effects that the evidence documents: it renews the sense of continuity with one's past self and of connection to those with whom the past was shared \citep{SedikidesEtAl2008}. These are reasons to value it, not conditions of its rationality, but they answer the charge that it is idle.
\end{reply}

\section{Conclusion}

Nostalgia is rational when it tracks real goods of the past that the subject can no longer have, and when it is in proportion to them. The goods are real and fixed, and what is lost is access to them from a later time. Its failures, distortion and excess, are failures of fit and of proportion, and they are not part of what nostalgia is.

\begin{thebibliography}{99}\small
\bibitem[Boym(2001)]{Boym2001} Boym, S. (2001). \emph{The Future of Nostalgia}. Basic Books.
\bibitem[D'Arms and Jacobson(2000)]{DArmsJacobson2000} D'Arms, J. and Jacobson, D. (2000). The moralistic fallacy: On the \enquote{appropriateness} of emotions. \emph{Philosophy and Phenomenological Research}, 61(1), 65--90.
\bibitem[de Sousa(1987)]{deSousa1987} de Sousa, R. (1987). \emph{The Rationality of Emotion}. MIT Press.
\bibitem[Howard(2012)]{Howard2012} Howard, S.~A. (2012). Nostalgia. \emph{Analysis}, 72(4), 641--650.
\bibitem[Sedikides et~al.(2008)]{SedikidesEtAl2008} Sedikides, C., Wildschut, T., Arndt, J. and Routledge, C. (2008). Nostalgia: Past, present, and future. \emph{Current Directions in Psychological Science}, 17(5), 304--307.
\bibitem[Wildschut et~al.(2006)]{WildschutEtAl2006} Wildschut, T., Sedikides, C., Arndt, J. and Routledge, C. (2006). Nostalgia: Content, triggers, functions. \emph{Journal of Personality and Social Psychology}, 91(5), 975--993.
\end{thebibliography}

\end{document}
